\documentclass{article}
\usepackage[utf8]{inputenc}
\usepackage{amsmath}
\usepackage{amssymb}
\usepackage{amsfonts}
\usepackage{geometry}
\geometry{a4paper, margin=1in}

\title{核心可解数学模型}
\author{Paper 10}
\date{\today}

\usepackage[UTF8]{ctex}
\begin{document}

\maketitle

\section*{保留的核心模型}
根据\texttt{}公式以及PDF文件中的描述，底层应用模型是第13页上的动态设施定位与处理预算分配模型。不过，论文指出，这种非凸双线性模型仅适用于小型实例，实际上难以求解。因此，所采用的\emph{模型实际上是可以求解的}核心模型。该模型基于附录/第27-29页的公式构建，而非原始的非凸模型。该模型保留了所有动态流行病约束和运营约束，并将双线性项替换为标准的精确/放宽线性重构形式。

\paragraph{决策变量。}
$x_i \in \mathbb Z_+$表示县$i$内的治疗设施数量，  
$d_{ik}\ge 0$表示在时间$k$时分配给县$i$的治疗预算，  
$h \ge 0$是超出预算的变量，  
$z_i$是$1/x_i$的换算系数，用于线性化公式中，  
$p_{ij},o_{ij},r_{ik}$是一元展开和McCormick松弛方法的辅助变量，  
$\mathbf z_k=(S_k,P_k,I_k,A_k,R_k,D_k)$是流行病状态向量。

\paragraph{集合与参数。}
$\mathcal C$是县级设置，而$\mathcal K=\{0,\dots,K\}$是时间段设置。  
权重$\lambda_D,\lambda_A,\lambda_{\mathrm{pr}},\lambda_{\mathrm{SVI}},
\lambda_{\mathrm{Pop}},\lambda_{\mathrm{inf}}$定义了目标之间的权衡关系。  
$N$是设施上限，$n_i$是初始设施数量，$d_k$是时间$k$时的状态处理预算，$d_{\min}$是每个设施的最低预算，$\theta$是允许的最大跨时期预算减少幅度，$d$是每患者的季度MAT成本，而$lb_{ik},ub_{ik}$则是预算的界限。

\paragraph{初始条件。}
这六个非负初始 compartments 的值都是给定的数据。它们代表了初始的种群状态。
\[
(S_0,P_0,I_0,A_0,R_0,D_0)
 = (S^{\mathrm{init}},P^{\mathrm{init}},I^{\mathrm{init}},
A^{\mathrm{init}},R^{\mathrm{init}},D^{\mathrm{init}}).
\]
\section*{保留以下实际计算公式：}
\begin{equation}
\min \quad
\lambda_D D_K
+ \lambda_A \sum_{k=0}^{K} A_k
+ \lambda_{\mathrm{pr}} \mathbf c_{\mathrm{pr}}^\top \mathbf x
+ \lambda_{\mathrm{SVI}} \mathbf c_{\mathrm{SVI}}^\top \mathbf x
+ \lambda_{\mathrm{Pop}} \sum_{k=0}^{K} \mathbf c_{\mathrm{Pop}}^\top \bar{\mathbf d}_k
+ \lambda_{\mathrm{inf}} h
\tag{Prac-Obj}
\end{equation}

\begin{equation}
\sum_{i=1}^{C} x_i - N \le h,\qquad h \ge 0
\tag{27}
\end{equation}

\begin{equation}
x_i \ge n_i
\quad \forall i \in \mathcal C
\tag{28a}
\end{equation}

\begin{equation}
\sum_{i=1}^{C} \bar d_{ik} \le d_k
\quad \forall k \in \mathcal K
\tag{28c}
\end{equation}

\begin{equation}
r_{ik} \ge d_{\min}
\quad \forall i \in \mathcal C,\ \forall k \in \mathcal K
\tag{28d}
\end{equation}

\begin{equation}
\bar d_{i,k-1} - \bar d_{ik} \le \theta
\quad \forall i \in \mathcal C,\ \forall k = 1,\dots,K
\tag{28e}
\end{equation}

\begin{equation}
lb_{ik} \le \bar d_{ik} \le ub_{ik}
\quad \forall i \in \mathcal C,\ \forall k \in \mathcal K
\tag{28f}
\end{equation}

\begin{equation}
S_{k+1} = S_k + (\epsilon P_k + \delta R_k - \alpha S_k)\Delta
\quad k=0,\dots,K-1
\tag{28g}
\end{equation}

\begin{equation}
P_{k+1} = P_k + (\alpha S_k - (\epsilon+\gamma+\beta)P_k)\Delta
\quad k=0,\dots,K-1
\tag{28h}
\end{equation}

\begin{equation}
I_{k+1} = I_k + (\beta P_k - \phi I_k)\Delta
\quad k=0,\dots,K-1
\tag{28i}
\end{equation}

\begin{equation}
A_{k+1} = A_k + \left(
\gamma P_k + \sigma R_k + \phi I_k - \zeta A_k
- \frac{1}{d}\sum_{i=1}^{C}\left(\bar d_{ik} - n_i r_{ik}\right)
- \mu A_k
\right)\Delta
\quad k=0,\dots,K-1
\tag{28j}
\end{equation}

\begin{equation}
R_{k+1} = R_k + \left(
\zeta A_k + \frac{1}{d}\sum_{i=1}^{C}\left(\bar d_{ik} - n_i r_{ik}\right)
- (\delta+\sigma)R_k
\right)\Delta
\quad k=0,\dots,K-1
\tag{28k}
\end{equation}

\begin{equation}
D_{k+1} = D_k + (\mu A_k)\Delta
\quad k=0,\dots,K-1
\tag{28l}
\end{equation}

\begin{equation}
x_i \in \mathbb Z,\qquad x_i \in [1,N]
\quad \forall i \in \mathcal C
\tag{28m}
\end{equation}

\begin{equation}
\bar d_{ik},r_{ik},S_k,P_k,I_k,A_k,R_k,D_k \ge 0
\quad \forall i \in \mathcal C,\ \forall k \in \mathcal K
\tag{28n}
\end{equation}

\section*{一元展开与McCormick约束}
\begin{equation}
x_i = \sum_{j=1}^{N} j\,p_{ij},\qquad
\sum_{j=1}^{N} p_{ij} \le 1,\qquad
\sum_{j=1}^{N} j\,o_{ij} = 1
\quad \forall i \in \mathcal C
\tag{29a}
\end{equation}

\begin{equation}
o_{ij} \ge \frac{1}{N}p_{ij}
\quad \forall i \in \mathcal C,\ \forall j=1,\dots,N
\tag{29b}
\end{equation}

\begin{equation}
o_{ij} \ge p_{ij} + z_i - 1
\quad \forall i \in \mathcal C,\ \forall j=1,\dots,N
\tag{29c}
\end{equation}

\begin{equation}
o_{ij} \le p_{ij}
\quad \forall i \in \mathcal C,\ \forall j=1,\dots,N
\tag{29d}
\end{equation}

\begin{equation}
o_{ij} \le z_i + \frac{1}{N}p_{ij} - \frac{1}{N}
\quad \forall i \in \mathcal C,\ \forall j=1,\dots,N
\tag{29e}
\end{equation}

\begin{equation}
r_{ik} \ge lb_{ik} z_i + \frac{1}{N}\bar d_{ik} - \frac{1}{N}lb_{ik}
\quad \forall i \in \mathcal C,\ \forall k \in \mathcal K
\tag{29f}
\end{equation}

\begin{equation}
r_{ik} \ge ub_{ik} z_i + \bar d_{ik} - ub_{ik}
\quad \forall i \in \mathcal C,\ \forall k \in \mathcal K
\tag{29g}
\end{equation}

\begin{equation}
r_{ik} \le ub_{ik} z_i + \frac{1}{N}\bar d_{ik} - \frac{1}{N}ub_{ik}
\quad \forall i \in \mathcal C,\ \forall k \in \mathcal K
\tag{29h}
\end{equation}

\begin{equation}
r_{ik} \le \bar d_{ik} + lb_{ik} z_i - lb_{ik}
\quad \forall i \in \mathcal C,\ \forall k \in \mathcal K
\tag{29i}
\end{equation}

\begin{equation}
\frac{1}{N} \le z_i \le 1,\qquad
p_{ij}\in\{0,1\},\qquad
0 \le o_{ij} \le 1
\quad \forall i \in \mathcal C,\ \forall j=1,\dots,N
\tag{29j}
\end{equation}

\section*{对保留的松弛度进行解释：}
这四个不等式（29f）～（29i）构成了$r_{ik}=\bar d_{ik}z_i$的McCormick松弛形式。它们定义了乘积松弛。在（28j）～（28k）中使用的处理项，就是所声明的松弛表达式$d^{-1}\sum_i(\bar d_{ik}-n_i r_{ik})$。解和最优值指的是在固定的初始条件下，针对该MILP松弛形式所得到的解和最优值。
\end{document}
